% 第 5 章 与现有工作对比及应用前景
% 本章引用键（实际使用）：tvla, dudect, rtlf, ctgrind, microwalk, kyberslash, ctverif, binsec, clangover, azouaoui2021, jancar2022
% 图表按章编号：本章为表 5.1；本章不放图

\section{与现有工作对比及应用前景}

\subsection{泄漏检测与评估方法对比}

软件时间侧信道的评估分成三条路线：统计检测、动态分析与形式化验证。本作品没有发明新的检测方法，贡献在把三条路线组织进同一个流程：所有结论都挂在编译产物上，每一步都能审计。第 3 章讲这套流程怎么搭，本章讲三条路线各自回答什么问题、答案能用在哪一步。

统计检测把程序当黑盒。TVLA\cite{tvla} 比较固定类与随机类的轨迹分布，dudect\cite{dudect} 把同一套判据用在在线计时上，用最大 $t$ 统计量与固定阈值给出“有/无”两种结论，RTLF\cite{rtlf} 比较分布与分位，并用 bootstrap 控制第一类错误。三者都回答“两类样本能不能分开”；差异来自哪条指令由定位工具回答。类怎么配对、样本按什么顺序跑、两次运行之间有没有漂移，都会改变结论。第 4.12 节的 M2 展示了不同判据的区别：RTLF 在未修复版与已修复版均检出分布差异，dudect 的阳性落在已修复版，预注册均值检验却对两个版本都未拒绝。

动态分析进到程序内部找秘密依赖。ctgrind 沿用 Valgrind 的秘密标记，报告秘密在什么时候进入分支条件或访存下标\cite{ctgrind}。Microwalk 用执行轨迹与互信息，找出与秘密关联最强的指令\cite{microwalk}。varlat 检查除法这类可变时延指令的操作数里有没有秘密\cite{kyberslash}。这类工具能指出问题在哪个函数、甚至哪条指令；覆盖范围限于实际执行过的输入，“这台机器上能不能测到”由计时层回答。定位阳性与计时阳性分属两类证据，M2 的三处除法正是两者分离的例子（第 4.4 节、第 4.7 节）。

形式化验证在给定模型与边界内给出另一种结论：证明控制流与访存不依赖秘密。ct-verif 把“常数时间”翻译成程序安全性的命题，在中间表示层完成证明\cite{ctverif}。Binsec/Rel 用关系符号执行，在二进制层做有界检查，并给出反例路径\cite{binsec}。这类方法成本高，模型范围通常不含指令时延、缓存与流水线行为；“是否依赖秘密”由证明回答，“本机计时是否可区分”由计时层回答。第 4.4 节 M4 的不安全判定与常数时间实现的有界安全、M5 的 1088 条路径全部通过（1088/1088），都属于有界结论。

表~\ref{tab:compare} 把三条路线与本作品的定位放在同一组维度上。每一行的最后一列写出这个差异导致的实践后果，而不只是一个勾。

\begin{table}[H]
\centering
\caption{泄漏检测与评估方法对比}
\label{tab:compare}
{\footnotesize
\begin{tabularx}{\textwidth}{|>{\raggedright\arraybackslash}p{2.7cm}|>{\raggedright\arraybackslash}p{2.3cm}|>{\raggedright\arraybackslash}p{2.3cm}|>{\raggedright\arraybackslash}p{2.5cm}|>{\raggedright\arraybackslash}p{2.3cm}|>{\raggedright\arraybackslash}X|}
\hline
方法 & 本机可测性 & 函数级定位 & 产物级证明 & 未检出的量化 & 多机制与复测 \\
\hline
TVLA\cite{tvla}、dudect\cite{dudect} & 是（黑盒计时） & 否 & 否 & 一般不报，阴性无预算 & 需另行组织。后果：换产物要整体重判 \\
\hline
RTLF\cite{rtlf} & 是（第一类错误有界） & 否 & 否 & 随样本量提高功效，不报 MDE & 需另行组织。后果：分布差异无法归因到指令 \\
\hline
ctgrind\cite{ctgrind}、Microwalk\cite{microwalk} & 否（不回答本机可测性） & 是（分支、访存、轨迹） & 否 & 不适用（只报依赖） & 单一机制类。后果：定位阳性易被读成可利用 \\
\hline
varlat\cite{kyberslash} & 否 & 是（除法等可变时延指令） & 否 & 不适用 & 针对单一指令类。后果：只覆盖已知机制，漏掉新形态 \\
\hline
ct-verif\cite{ctverif}、Binsec/Rel\cite{binsec} & 否 & 是（函数或内核） & 是（给定边界内。ct-verif 在中间表示层） & 不适用（二值结论） & 单一函数或内核。后果：不建模指令时延，答不了本机计时 \\
\hline
本作品 & 是（预注册 + Holm + MDE） & 是（多工具交叉） & 部分（M4、M5 有界） & 报告 MDE 与样本量 & 5 模块统一规程、防护前后同规程复测。结论绑定产物并按证据层追溯 \\
\hline
\end{tabularx}
}
\end{table}

五列维度各有用途。本机可测性决定结论能不能用来做发布判断，函数级定位决定阳性结论能不能落到具体代码，产物级证明决定结论是否覆盖所有输入，未检出的量化决定阴性结论有多少说服力，多机制与防护复测决定一次评估能不能复用。五列合起来，读者才能判断一条结论该信到哪一步。

本作品把三条路线当作组件来用。同一个二进制既进定位工具，也进计时窗口。定位结果、计时判定与防护复测按证据层分开登记。每个阴性附最小可检测效应，每个未执行单列。再用 90 格排查与 9×5 工具矩阵做交叉（第 4.5 节、第 4.12 节）。代价是采集和跑工具要多花时间，收益是每个结论都能说清“在哪一层、对哪个产物、在什么预算下成立”。

与单一工具相比，本平台的核心优势在于多维度证据的协同印证与全链路可追溯性：通过将同一份二进制产物的反汇编特征、动态污点轨迹、有界形式化验证与锁频统计计时绑定归档，平台有效克服了单一统计工具无法定位代码、单一动态工具无法评估本机实际时延的局限。

与单机制研究的分界在问题范围。KyberSlash 研究秘密除法，clangover 追踪编译器引入的分支，Azouaoui 等量化 Kyber 解密与重加密各步骤的泄漏\cite{kyberslash,clangover,azouaoui2021}；本作品侧重检测方法的可审计性与多机制覆盖。功耗、电磁与故障注入采用不同的观测信号，属于互补方向。

\textbf{机制定位与计时结果分别报告。}M2 的三处除法在定位层稳定出现，varlat 报了 1280 处错误，集中在 3 个调用上下文；端到端与除法热点微基准均未检出，最小可检测效应为 24.7--67.1 ns。% A-18 % A-03 % A-04

\textbf{锁频提供统一且可核对的工作点。}不锁频与锁频的变异系数（CV）比为 0.964、1.006、1.089，三个比值没有统一的走向。锁频提供的是可核对、可恢复的工作点，让“这次采集在什么频率下完成”可被审计。% A-29

\textbf{同一机制在不同优化级别与架构上形态不同。}M1 在 -Os 下被编译成按位测试后条件跳转，-O2/-O3 下定位工具的命中更多（8 个、15 个）却计时未检出。aarch64 上位域提取与条件选择取代了分支，机制不复现。分支更多不一定更容易检出，结论必须绑定产物与环境。% A-26 % A-27

表~\ref{tab:compare} 展示三条路线各自的问题域，也呈现本作品如何将统计检测、机制定位与有界验证的结果放在同一证据框架中。“部分”指 M4、M5 的有界验证范围，不影响赛题 12 条要求的覆盖状态；附录 A 对照表中的 12 项均为满足。各路线的测量条件与方法覆盖不同，表中按证据范围比较，不据此排列总体优劣。

\subsection{应用场景与推广边界}

本作品可以作为密码库发布前检查、构建配置回归、防护方案评估与教学评测的本地工具。它不替代独立审计或形式化认证。四个场景的共同前提是：使用者手里有源码或二进制，能在一台受控机器上重复运行，并愿意记录编译条件与样本量。对 44 名密码库开发者的调查显示，开发者认为时间侧信道的防护并不难实施，真正的差距在工具有没有被采用、工程投入够不够\cite{jancar2022}。

后量子密码库发布前检查由库开发者使用。时机是版本冻结、发布候选版编译完成之后。输入是待发布的二进制、编译条件与公开向量（官方给出的固定测试输入）。输出是定位层的机制清单，以及固定机器上的计时判定。定位层用反汇编、Memcheck 与 varlat，成本低，可以在每次候选版上跑。计时层需要锁频与裸金属，按版本或按周在固定机器上跑。

小团队可以只跑定位层，把计时层留给有裸金属的机器。两层的结果分开登记，分别回答目标机制是否存在、本机计时能否区分。第 4.12 节的 9×5 矩阵说明，单工具既可能漏掉机制，也可能给出相反的读法，发布判定至少要两类工具结论一致。

编译器与构建配置变更回归由平台维护者使用。时机是升级编译器、改优化级别或换目标架构之前。clangover 表明这些变更会静默改变产物中的分支形态\cite{clangover}。本作品按“二进制摘要 + 优化级别”组织判定，先比对新旧产物的反汇编与 Memcheck，再在同一锁频工作点复测计时。得到的是新旧产物的差异，得不到“新版本安全”的结论：新产物可能引入新机制，也可能只是把旧机制挪了位置。面向持续集成（CI）的自动回归闸门还没有实现，列为后续工作。

防护方案的工程评估由实现者与评审者使用。时机是候选补丁进入主线之前。输入是同一份源码的防护前后两版、同一组公开向量与同一锁频规程。输出是定位结果、计时判定与成对开销。M1 的掩码选择使解封装 +1.51\%，M2 的除法替代使 $-2.92\%$，前者变慢、后者变快，说明防护的代价方向因实现而异，按实现独立记录。评估输出覆盖“这个方案在这个产物上的代价与效果”；方案之间的最优选择比较按实现者决策处理。% A-25

教学与评测把 M3、M4 的人工校准夹具与 M5 阴性对照当作工具测试集。时机是工具版本冻结或课程实验开始之前。教师或工具开发者用已知机制的夹具检查工具能否识别，用阴性对照检查会否误报。输入是夹具源码与固定输入类，输出是识别、漏报与误报的清单。M5 上运行过的 7 种工具误报为 0。% A-24

各场景的测量环境、攻击面与未验证项集中见第 1.2 节和第 6.2 节。本作品围绕 12 条赛题要求完成三层防护闭环：真实历史缺陷 M1、M2 接入标准补丁，并进行二进制产物级消融核验；人工校准夹具 M3、M4 实施自主常数时间改造并复测，其中 M4 另完成有界验证；平台自身的 1,048 个调用点接受静态污点审计，发现并修复 1 处退出码侧信道。对应状态与证据见附录 A（表 A1），12 条要求全部覆盖并满足。KyberSlash 产物中同源消息编码分支是 90 格全流程排查得到的超额发现，说明除法补丁之外仍须执行完整产物核验；该分支按已登记发现处理，防护项状态以 90 格核验结果为准。% A-18 % A-23 % A-36

最终交付是一组按产物、证据层与适用边界组织的可追溯本地记录。
